\documentclass{article}
\usepackage[UTF8]{ctex}
\usepackage{amsmath,amssymb,amsfonts}
\usepackage{enumitem}
\usepackage{geometry}
\geometry{a4paper, margin=1in}
% === 错题本排版优化（不参与题目标识） ===
\setlength{\parindent}{0pt}
\setlength{\parskip}{0.35em plus 0.10em minus 0.08em}
\setlength{\emergencystretch}{3em}
\setlist{itemsep=0.25em, topsep=0.40em, parsep=0pt}
\allowdisplaybreaks[2]
% === 错题本排版优化结束 ===
\begin{document}
\section*{西北工业大学2025年数学分析真题}
\section*{182．西北工业大学2025年数学分析试卷}

\subsection{题目1}
1．（20 分）用极限的严格数学定义证明：

\begin{itemize}
  \item[（1）] \(\lim _{x \rightarrow-\infty} \ln \left(1+\frac{1}{x}\right)=0\) ．
  \item[（2）] \(\lim _{n \rightarrow \infty} \frac{e^{n}}{n!}=0\) ．
\end{itemize}

\subsection{题目2}
2．（20 分）解答如下问题：

\begin{itemize}
  \item[（1）] 求极限 \(\lim _{x \rightarrow 0^{+}} \frac{x^{x}-x^{\sin x}}{x^{2}}\) ．
  \item[（2）] 求级数的和 \(\sum_{n=2}^{\infty} \frac{1}{n(n-1) 3^{n}}\) ．
\end{itemize}

\subsection{题目3}
3．（15 分）试用确界原理证明有限覆盖定理．

\subsection{题目4}
4．（15 分）证明：若 \(f\) 是区间 \([a, b]\) 上只有有限个间断点的有界函数，则 \(f\) 在 \([a, b]\) 上黎曼可积．

\subsection{题目5}
5．（15 分）计算 \(\int_{0}^{+\infty} \frac{\cos a x-\cos b x}{x^{2}} \mathrm{~d} x\) ，其中 \(b>a>0\) ．

\subsection{题目6}
6．（15 分）设 \(n\) 为正整数，\(x, y>0\) ，用条件极值方法证明 \(\frac{x^{n}+y^{n}}{2} \geq\left(\frac{x+y}{2}\right)^{n}\) ．

\subsection{题目7}
7．（15 分）设 \(f(x)\) 在 \([0,+\infty)\) 上一致连续，证明：存在 \(a, b>0\) 使得 \(|f(x)| \leq a|x|+b\) ．

\subsection{题目8}
8．（15 分）计算曲面积分

\[
I=\iint_{S} \frac{x \mathrm{~d} y \mathrm{~d} z+y \mathrm{~d} z \mathrm{~d} x+z \mathrm{~d} x \mathrm{~d} y}{\sqrt{\left(x^{2}+y^{2}+z^{2}\right)^{3}}} .
\]

其中 \(S\) 为曲面 \(\frac{x^{2}}{2}+\frac{y^{2}}{3}+\frac{z^{2}}{4}=1\) ，取其外侧。

\subsection{题目9}
9．（15 分）设 \(f(x)\) 是定义在 \((-\infty,+\infty)\) 上的凸函数且有上界，证明：\(f(x)\) 是常数．

\subsection{题目10}
10．（5 分）设 \(f(x)\) 在 \([a, b]\) 上连续，非负且严格单调递增，由积分中值定理，对任意的正整数 \(k\) ，存在\(x_{k} \in[a, b]\) ，使得

\[
f^{k}\left(x_{k}\right)=\frac{1}{b-a} \int_{a}^{b} f^{k}(x) \mathrm{d} x
\]

证明： \(\lim _{k \rightarrow \infty} x_{k}=b\) ．
\end{document}
